\documentclass[fontset=none]{ctexart}
\usepackage{geometry,graphicx,fontspec,lastpage,hyperref}
\geometry{a4paper,left=3.18cm,right=3.18cm,top=2.54cm,bottom=2.54cm,headheight=28pt}
\ctexset{section={
	name={第,章},
	number=\chinese{section},
	},
	subsection={
	number=\chinese{subsection},
	},
	fontset = adobe,
	today = big
}

\hypersetup{
	hidelinks,
	bookmarksnumbered = true,
	pdfauthor= {Ouyang Song},
	pdftitle = {诉讼策略分析报告}
}
\usepackage{fancyhdr}
\fancyhf{}
\lhead{\includegraphics[height=0.3cm]{img/logo.png} \ 湖北山河律师事务所}
\chead{诉讼策略分析报告}
\rhead{ \thepage / \pageref{LastPage}}
\lfoot{ \scriptsize 中国武汉市江汉区新华路385号南达大楼11-13层\\
11-13/F, Nanda Building, No.385 Xinhua Road, Jianghan District, Wuhan China\\
电话（Tel）：+86-27-59516866　59516867　  传真（Fax）：+86-27-59516823}
\renewcommand{\headrulewidth}{0.8pt}
\pagestyle{fancy} % \sectionmark 的重定义需要在 \pagestyle 之后生效
\renewcommand\sectionmark[1]{%
	\markright{\CTEXifname{\CTEXthesection——}{}#1}}
	\renewcommand{\footrulewidth}{0.6pt}

\author{欧阳松}
\date{2018-8-4}
\title{报告}

\begin{document}

\fontspec{Times New Roman}

\begin{titlepage}
\begin{center}
	\vspace*{3ex}
	\includegraphics[width=.5\linewidth]{img/logo-within.png}
	\vspace{3ex}

	\noindent\rule{.6\textwidth}{1.6pt}\vspace*{-\baselineskip}\vspace*{2pt}\\ % Thick horizontal rule
	\rule{.6\textwidth}{0.4pt} % Thin horizontal rule

	\vspace{2ex} % Whitespace above the title

	{\zihao{2} \heiti 诉讼策略报告}

	\vspace{1ex} % Whitespace below the title

	\noindent\rule{.6\textwidth}{0.4pt}\vspace*{-\baselineskip}\vspace{3.2pt}\\ % Thin horizontal rule
	\rule{.6\textwidth}{1.6pt} % Thick horizontal rule
	
	\vspace{3em}
	\makebox[5em][s]{客户名称：}\makebox[.5\linewidth][l]{湖北省科技投资集团有限公司}\\
	\makebox[5em][s]{}\makebox[.5\linewidth][l]{武汉光谷资产投资管理有限公司}
	\makebox[5em][s]{案由：}\makebox[.5\linewidth][l]{房屋租赁合同纠纷}
	\makebox[5em][s]{团队成员：}\makebox[.5\linewidth][l]{xlkjslkjdflskjfs、sakjgkasgkahg、kslgkashgkashgkah}
	\makebox[5em][s]{联系人：}\makebox[.5\linewidth][l]{欧阳松}
	\makebox[5em][s]{电子邮箱：}\makebox[.5\linewidth][l]{ouyang.sh@foxmail.com}

	\vspace{5em}


	{湖北山河律师事务所}\\
	\today
	\end{center}

	\vspace{13ex}
	\makebox[.5\linewidth][l]{地址：武汉市江汉区新华路385号南达大楼11-13楼}\par
	\makebox[.5\linewidth][l]{Address:11-13/F, Nanda Building, No.385 Xinhua Road, Jianghan District, Wuhan China}\par
	\makebox[.5\linewidth][l]{Tel:86-027-59516866 \hfill Fax:86-027-59521682}




\end{titlepage}

\newpage

\pagestyle{empty}
\setcounter{page}{0}
\tableofcontents

\newpage

\pagestyle{fancy}
\noindent \makebox[\linewidth][c]{\zihao{4}关于A公司与B公司房屋租赁合同纠纷一案}

\vspace{1em}

\noindent \makebox[\linewidth][c]{\zihao{4}之}

\vspace{1em}
\noindent \makebox[\linewidth][c]{\textbf{\zihao{3} 诉讼策略分析报告}}\par
\vspace{1.5em}

\noindent 致：湖北省科技投资集团有限公司、\par
武汉光谷资产投资管理有限公司
\vspace{1em}



湖北山河律师事务所（简称“本所”）作为湖北省科技投资集团有限公司（简称“省科投集团”）及武汉光谷资产投资管理有限公司（简称“光谷资产公司”）的法律顾问，现就省科投集团与武汉银瑞丰商贸有限公司（曾用名“武汉银瑞丰投资有限公司”，简称“银瑞丰公司”）房屋租赁合同纠纷一案，出具如下初步诉讼策略分析，以供贵司决策时参考使用。



\section{制定依据} % (fold)
\label{sec:制定依据}

	\subsection{事实依据} % (fold)
	\label{sub:事实依据}
	为制定本诉讼策略分析报告，本所律师仅审查了贵司提供的下列文件资料（均为扫描件）:
	\begin{enumerate}
	\item 档号为4-2010-28的《合同审批单》；
	\item 2010年10月22日，省科投集团与武汉鼎中创智投资有限公司（简称“鼎中创智公司”）签署的合同编号为科投租字（2010）-018号《房屋租赁合同》；

	\item 2010年10月28日，省科投集团向鼎中创智公司发出的《联系函》；

	\item 档号为4-2012-45的《湖北省科技投资集团有限公司合同审批单》；

	\item 2012年9月16日，省科投集团与鼎中创智公司、银瑞丰公司签署的编号为科投租字（2012）-32号《补充协议》；

	\item 2015年7月22日，省科投集团资产部向银瑞丰公司发出的《催收函》（载有罗婧于2015年7月27日的签名）；

	\item 2016年11月29日，省科投集团向银瑞丰公司发出的《工作联系函》（载有罗婧的签名）；

	\item 2017年10月12日，光谷资产公司向银瑞丰公司发出的《催款通知书》（罗婧于2017年10月16日签收）；

	\item 2017年10月18日，银瑞丰公司向省科投集团作出的《承诺函》；

	\item 2018年5月14日，省科投集团及光谷资产公司向银瑞丰公司作出的《工作联系函》（银瑞丰公司于2018年5月16日盖章签收）。
	\end{enumerate}
	% subsection 事实依据 (end)

	\subsection{法律依据} % (fold)
	\label{sub:法律依据}
	
	\begin{enumerate}
		\item 《中华人民共和国民法总则》
		\item 《中华人民共和国合同法》
		\item 《中华人民共和国物权法》
		\item 《中华人民共和国民事诉讼法》
		\item 《最高人民法院关于适用〈中华人民共和国民事诉讼法〉的解释》
		\item 《最高人民法院关于适用〈中华人民共和国合同法〉若干问题的解释（二）》
		\item 《最高人民法院关于审理城镇房屋租赁合同纠纷案件具体应用法律若干问题的解释》
		\item 国务院《中华人民共和国诉讼费用交纳办法》
	\end{enumerate}
	% subsection 法律依据 (end)

	\subsection{其他依据} % (fold)
	\label{sub:其他依据}

	\begin{enumerate}
		\item 贵司工作人员陈述
		\item 本所律师查询到的工商信息
		\begin{itemize}
			\item 企业报告-A公司
			\item 企业报告-B公司
		\end{itemize}
	\end{enumerate}
	
	% subsection 其他依据 (end)
% section 制定依据 (end)

	\section{案件事实} % (fold)
	\label{sec:案件事实}
	
	% section 案件事实 (end)

	\section{案件法律分析} % (fold)
	\label{sec:案件法律分析}
	
	% section 案件法律分析 (end)

	\section{诉讼费用预算} % (fold)
	\label{sec:诉讼费用预算}
	
	% section 诉讼费用预算 (end)

	\section{特别说明} % (fold)
	\label{sec:特别说明}
	1、需要提醒贵司的是，本诉讼策略报告所涉相关事实的了解与判断，完全依赖于贵司向本所提供的文件、材料及贵司工作人员所作陈述与说明。本所对贵司未能提供的其他文件、材料的缺失而导致本所法律分析不准确、不完整或不全面不承担相关法律责任；贵司应对其提供的文件、材料及相关人员陈述、说明的真实性负责。

	2、本诉讼策略分析报告系本案的初步法律分析意见，针对不同的情况下，对方不同诉讼策略、证据的提交等，需贵司配合核查事实情况后，本所及时有效的调整后续方案，以确保目标的实现。

	3、本所律师对本诉讼策略分析报告享有著作权，本诉讼策略分析报告是在充分信任贵司的基础上作出的未提及保留意见的法律分析及报价，仅供贵司在处理与武汉银瑞丰商贸有限公司房屋租赁合同纠纷一案相关事宜时使用，未经本所及本律师同意，不得向任何第三人提供，亦不得作为其他任何用途使用。


	% section 特别说明 (end)


	\vspace*{5em}
	
	\hfill \makebox[10em][c]{湖北山河律师事务所}\\
	\vspace{1.5ex}
	\hfill \makebox[10em][c]{\today}
	
\end{document}